\section{Related Work}\label{sec:related}

\paragraph{Sparse autoencoders for interpretability.} Sparse dictionary learning on transformer activations was shown
to yield many interpretable latents by \citet{bricken2023monosemanticity} and \citet{cunningham2023sparse}. Later work
changed the activation function to improve the trade-off between reconstruction and sparsity: TopK SAEs fix the number
of active latents \citep{gao2024scaling}, BatchTopK applies the constraint over a batch \citep{bussmann2024batchtopk}, and
JumpReLU SAEs learn a per-latent threshold with a straight-through estimator \citep{rajamanoharan2024jumping}. Public
suites exist for GPT-2 Small \citep{bloom2024gpt2} through SAELens \citep{saelens2024}, and for Gemma 2 through Gemma
Scope \citep{gemmascope2024}. Our in-house JumpReLU uses a simpler sigmoid surrogate, and we use the public suites as
external references rather than as objects of improvement.

\paragraph{Scaling and evaluating SAEs.} \citet{gao2024scaling} study how SAE quality scales with width and training
compute and argue that reconstruction and sparsity should be complemented by downstream measures, including the
change in the model's loss when the reconstruction is spliced in. Our study shares that motivation and applies it to a
separation of four faithfulness levels with explicit controls. SAEBench \citep{karvonen2025saebench} evaluates SAEs on
a broad suite of downstream metrics and finds that reconstruction and sparsity proxies do not reliably predict practical
usefulness, and AxBench \citep{wu2025axbench} finds that simple baselines outperform SAE-based methods for concept
detection and steering. \citet{engels2024darkmatter} analyse the structure of the SAE reconstruction error, and
\citet{korznikov2026random} propose random-decoder sanity baselines, which we include in the scaling experiment.
\citet{chaudhary2024evaluating} compare open-source GPT-2 SAEs with supervised subspaces on disentangling factual
knowledge, which motivates our low-rank causal reference. Our study is narrower than these benchmarks, with two models
and one task family per question, but it pre-registers each comparison and pairs every faithfulness measure with a
noise, PCA, synonym or random-subspace control.

\paragraph{Training SAEs for downstream faithfulness.} \citet{braun2024e2e} train SAEs end-to-end with a KL objective on
the model's output distribution, and \citet{karvonen2025kl} show that a short fine-tune with such an objective recovers
much of its benefit. Our output-KL and R4 fine-tunes follow this idea but are not reproductions; their contribution here
is to measure what a behavioural objective does to the other three faithfulness levels when sparsity is held fixed.

\paragraph{Causal interpretability.} Activation patching locates where a model carries information by copying
activations between runs, as in causal tracing of factual associations \citep{meng2022locating}; attribution patching
approximates such interventions efficiently \citep{syed2023attribution}, and sparse feature circuits use SAE latents as
the units of a causal graph \citep{marks2025sparse}. Our causal-site sweep follows the patching paradigm, and our
decomposition asks a prior question for SAE-based circuit work: how much of a located effect the SAE reconstruction
retains in the first place.
% TODO(citation): add a verified reference for activation patching methodology (e.g. causal mediation / interchange
% intervention work) once checked against the original paper; none is in the repository bibliography yet.

\paragraph{Probing and representation engineering.} \citet{marks2023geometry} report linear truth structure in language
model activations using the Geometry of Truth datasets that we use, and \citet{zou2023representation} introduce
representation engineering with persona-contrast prompts, from which our original persona task was derived. Our
distinction between frozen and re-trained readers relates to a broader literature on what probe accuracy does and does
not show about a representation.
% TODO(citation): add verified references on probing methodology (control tasks; probing-classifier surveys) once
% checked against the original papers; none is in the repository bibliography yet.
